\section{Open questions (5 truly-open, with concrete next steps)}
\label{sec:open}

\subsection*{Q1. Open-source MC re-derivation of Table 2 S-values}
\textbf{Question:} Can the Table 2 S-values be independently re-derived from an
open-source Monte Carlo pipeline (Geant4-DNA or TOPAS-nBio) using the published
PC3-PIP cellular geometry parameters (Table 1) rather than trusted verbatim
from the paper's proprietary Geant4 run?

\textbf{Basis:} The paper's Geant4 v10.03(6) run with Livermore EM +
FTFP\_BERT hadronic physics used custom cellular geometry meshes derived from
microscopy that are \emph{not} published (only mean radii/volumes are in
Table 1). Our replication trusts Table 2 verbatim, which is the documented
LUCID pattern where the proprietary MC substrate is never re-derived.
Independent re-derivation would close the single largest open computational
gap and let us quantify uncertainty on S-values themselves (currently
zero-variance in our pipeline).

\textbf{Next steps:} Build a Geant4-DNA cellular geometry using the paper's
Table 1 dimensions (spherical approximation initially, then ellipsoidal).
Emit Ac-225 and its four alpha daughters (Fr-221, At-217, Bi-213, Po-213)
plus Lu-177 beta spectrum. Run $10^6$ decay histories per source-target
compartment pair, score energy deposition in cytoplasm / nucleus / membrane /
medium targets. Compare recomputed S-values to Table 2 entries. Target:
reproduce $S(\text{Ac,cell})/S(\text{Lu,cell})$ ratio (200--550$\times$)
within 20\% and absolute cytoplasm S-values within 2$\times$. Free/open:
Geant4-DNA is BSD-licensed, runnable on a single node. Estimate 3--5 days
compute for one geometry pass.

\subsection*{Q2. Does the MIRD offset collapse under full TAC + MIRDcell?}
\textbf{Question:} Does the constant multiplicative MIRD offset
(Lu 1.28$\times$, Ac 2.4$\times$) collapse to unity when the full
time-activity curve (slow-ramp uptake during 3~h incubation + explicit
cell--cell cross-dose averaging via MIRDcell) is substituted for the
instant-uptake / self-dose approximation used in our replication?

\textbf{Basis:} Our replication finds a strictly constant multiplicative
offset across all activity concentrations, diagnostic of a single modeling
choice (instant maximum uptake) rather than an algebra error. The paper
explicitly performs MIRDcell cross-dose averaging over the cell-dimension
distribution. If the offset resolves cleanly by relaxing these two
assumptions, the MIRD pipeline can be up-tiered from
REPRODUCED-IN-STRUCTURE to REPRODUCED-QUANTITATIVELY.

\textbf{Next steps:} Implement a two-piece TAC: linear ramp from 0 to peak
uptake over the 3~h incubation window, then biexponential washout with
$t_{\text{bio}}$=2.3 h ($F$=0.59) plus $t_{\text{phys}}$. Recompute integrated
activity $\tilde{A}$ per concentration. Separately, download MIRDcell
(Vanderbilt, free) or reimplement its cell-dimension-distribution weighting
from Table 1 percentile ranges. Rerun both isotope pipelines. Target: Lu
offset from 1.28$\times$ to within 10\% of 1.0; Ac offset from 2.4$\times$ to
within 20\% of 1.0. If offset persists, the residual points to a real
modeling gap in the paper worth flagging.

\subsection*{Q3. LQ-model sensitivity of RBE}
\textbf{Question:} How sensitive is the RBE central estimate (paper 4.2, our
3.0--3.3) to the choice of survival-curve model (linear-exponential vs.
linear-quadratic $S(D) = \exp(-\alpha D - \beta D^2)$), and does modern
radiobiology suggest a $\beta$-term should be non-negligible for either
isotope at the tested dose range?

\textbf{Basis:} The paper (and our replication) fits pure linear-exponential
$S(D) = \exp(-\alpha D)$. Standard radiobiology (LQ model) allows a
quadratic term that becomes non-negligible for beta emitters at doses above
2--3 Gy where sublethal-damage repair saturates. The Lu-177 dose range in
Table 3 spans 0.74--3.70 Gy, which straddles this boundary. If $\beta > 0$
for Lu-177 while alpha-only holds for Ac-225 (theoretically justified by
the high LET of alpha making sublethal damage rare), the effective
$\alpha$(Lu) from a pure-linear fit is inflated and the true RBE is larger
than 3.0--3.3 and closer to (or above) the published 4.2.

\textbf{Next steps:} Refit Lu-177 survival with a two-parameter LQ model
using both digitization reads. If $\beta > 0$ with 95\% CI excluding zero,
report the $\alpha/\beta$ ratio and recompute RBE
$= \alpha_{\text{Ac}} / \alpha_{\text{Lu}}$ using the LQ-$\alpha$ for Lu.
Compare against published $\beta$-emitter $\alpha/\beta$ values in
prostate cancer cell lines (literature range: 1.5--10 Gy). For Ac-225,
fit LQ but expect $\beta$ to be consistent with zero given alpha-LET
saturation. Free tool: \texttt{scipy.optimize.curve\_fit}, already installed.

\subsection*{Q4. RBE portability from 2D monolayer to clinically-relevant geometries}
\textbf{Question:} Does the $\sim$4$\times$ RBE derived from clonogenic
survival in PC3-PIP monolayer culture translate to a comparable RBE in more
clinically relevant settings (3D spheroids, PSMA-expressing patient-derived
xenograft slices, or micrometastasis models where the alpha-particle range
$\sim$50--80~\textmu m approaches the tumor cluster diameter)?

\textbf{Basis:} The paper's RBE is measured in 2D monolayer culture where
every cell is directly PSMA-labeled and hit by the emitted radiation.
\emph{In vivo}, alpha-particle short range means the effective self-dose to
micrometastases depends on cluster size relative to the $\sim$5-cell alpha
range, while Lu-177's $\sim$2~mm beta range gives crossfire dose to
unlabeled bystander cells. A single-cell-monolayer RBE may over- or
under-state the clinically effective RBE by a factor of 2--3 depending on
geometry.

\textbf{Next steps:} Literature synthesis pass (no wet-lab): systematic
review of published Ac-225 vs. Lu-177 RBE measurements in (a) 2D monolayer,
(b) 3D spheroid, (c) xenograft, (d) patient-derived organoid systems for
prostate cancer or PSMA-expressing tumors. Extract RBE values with 95\% CIs
and geometry description. Test hypothesis: RBE(monolayer) $>$ RBE(spheroid)
$>$ RBE(xenograft) as bystander crossfire dilutes the alpha advantage.
Endpoints free via PubMed / Semantic Scholar API. Deliverable: forest plot
of published RBE with our replicated 3.0--3.3 anchored in the 2D-monolayer
bin.

\subsection*{Q5. Bayesian posterior on RBE marginalized over digitization noise}
\textbf{Question:} Can an ML-fit (Bayesian hierarchical or Gaussian-process
regression) over the digitized Figure 3 survival points, marginalized over
digitization noise, produce a posterior on RBE that is honestly narrower
than the two-read range (2.96--3.33) we currently report, and does it agree
with the published 4.2 within its 95\% credible interval?

\textbf{Basis:} Our replication reports RBE as a discrete range across two
independent digitization reads (essentially a poor-man's uncertainty
quantification). A principled Bayesian fit that treats each digitized point
as a noisy observation of an underlying survival fraction, marginalizes
over digitization systematics, and hierarchically pools information across
the six (Lu) or eight (Ac) concentrations should produce a posterior
distribution over $\alpha$ for each isotope, and thus a posterior on RBE.
This would let us honestly say whether the paper's 4.2 is inside or outside
our 95\% credible interval, which the current point-estimate framing cannot.

\textbf{Next steps:} Build a Stan or PyMC model: for each isotope,
$S_{ij} \sim \text{LogNormal}(\log(\exp(-\alpha_i D_i)), \sigma_{\text{dig}})$
with weakly informative priors on $\alpha_i$ and $\sigma_{\text{dig}}$. Fit
both isotopes jointly with a shared digitization-noise hyperprior. Sample
the posterior on RBE $= \alpha_{\text{Ac}} / \alpha_{\text{Lu}}$. Report the
95\% credible interval. Cross-check against a non-parametric bootstrap
over the two digitization reads. Free: PyMC on a laptop, $\sim$10~min wall
clock. Deliverable: RBE posterior figure with the paper's 4.2 $\pm$ 0.46
overlaid as a reference band.
